\documentclass[../../main.tex]{subfiles}
\graphicspath{{../../image/}{../../figure/}}
\begin{document}

\section{熵可容性条件}\label{section:ch04-熵可容性条件}

正如我们在第 III 章中看到的，连续物理中出现的每一个平衡律组都伴随一个熵不等式，任何有物理意义的过程都必须满足它，因为它显式或隐式地表达了热力学第二定律。这启示了下述刻画双曲守恒律组可容弱解的程序。

设方程组\eqref{align:ch04-守恒律组}配有一个指定熵 $\eta$ 及相应熵流 $Q$，使得\eqref{align:ch04-熵流关系}成立。一个\eqref{align:ch04-守恒律组}的、在第 4.3 小节意义下定义于 $\mathbb{R}^{m}\times[0,T)$ 上的弱解，称为相对于 $\eta$ 满足\emph{熵可容性判据}，如果在分布意义下于 $\mathbb{R}^{m}\times[0,T)$ 上成立
\begin{align}
\partial_{t}\eta(U(x,t))+\operatorname{div}Q(U(x,t))\leqslant 0.
\label{align:ch04-熵不等式}
\end{align}

显然，\eqref{align:ch04-守恒律组}的任一古典解都是可容的，因为它满足等式\eqref{align:ch04-熵守恒}。另一个相关注记是：熵可容性判据在解上诱导出一个时间\emph{不可逆}条件：若 $U(x,t)$ 是\eqref{align:ch04-守恒律组}满足\eqref{align:ch04-熵不等式}为严格不等式的可容弱解，则 $\bar{U}(x,t)=U(-x,-t)$ 虽然也是解，却不可容。

一个自然问题是：如何为可容性判据指定合适的熵。例如，显然相对于熵 $\eta$ 可容的弱解，相对于熵 $\bar{\eta}=-\eta$ 未必可容。当方程组来自物理时，应由物理指定自然熵。在缺乏物理指引时，则须使用数学论证。当然，期望由指定熵诱导的可容性判据与稍后引入的其他可容性条件相容。另一个自然条件是：可容弱解应具有合理的稳定性。正如我们将看到，当指定熵 $\eta(U)$ 为凸、或至少“类凸”时，以上所有要求都得到满足。

读者应牢记：凸性仅在方程组取正则形式时才是熵的相关性质。一般情形下，$\eta$ 的凸性应替换为条件：\cref{chapter:ch03-双曲平衡律方程组}中 $(n$ 维向量值$)$ 函数 $B(U,x,t)$ 的 $(n\times n$ 矩阵值$)$ 导数 $\mathrm{D}B(U,x,t)$ 正定。


回顾\cref{chapter:ch03-双曲平衡律方程组}中所考察的例子可知：作为把平衡律组化为正则形式的状态向量函数，在热弹性流体（例 3.3.5）、等温热弹性流体（例 3.3.6）与磁流体力学（例 3.3.9）的情形中，熵确实是凸的。这可能使人以为连续物理方程中熵一般是凸的。然而正如我们将看到，情况并非总是如此；因而有必要考虑更广的熵函数类。


对任一满足熵可容性判据的弱解 $U$，\eqref{align:ch04-熵不等式}左端是一个非正分布，从而是一个测度，称为\emph{熵产生测度}。于是\cref{chapter:ch01-平衡律}蕴含：映射 $t\mapsto\eta(U(\cdot,t))$ 在 $[0,T)\setminus\mathscr{F}$ 上在 $L^{\infty}(\mathbb{R}^{m})$ weak$^{*}$ 下连续，其中 $\mathscr{F}$ 至多可数。进而，对每一在 $\mathbb{R}^{m}\times[0,T)$ 内紧支集的非负 Lipschitz 检验函数 $\psi(x,t)$ 以及任一 $\tau\in[0,T)\setminus\mathscr{F}$，

\begin{align}
\int_{\tau}^{T}\int_{\mathbb{R}^{m}}\left[\partial_{t}\psi\eta(U)+\sum_{\alpha=1}^{m}\partial_{\alpha}\psi Q_{\alpha}(U)\right]\,\mathrm{d}x\mathrm{d}t+\int_{\mathbb{R}^{m}}\psi(x,\tau)\eta(U(x,\tau))\,\mathrm{d}x\geqslant 0.
\label{align:ch04-熵积分tau}
\end{align}

确定集合 $\mathscr{F}$ 何时实际为空将是重要的，但 Euler 方程所具有的剧烈振荡、奇异解（将在第 4.8 小节中展示）表明回答这一问题并不容易。事实上，在本书写作之时，仅在标量情形 $n=1$ 下有严格证明 $\mathscr{F}=\varnothing $（见\cref{chapter:ch06-标量守恒律L1理论}）。正如我们将看到，若至少 $0\notin\mathscr{F}$，则将大大有助于分析；此时取 $\tau=0$，\eqref{align:ch04-熵积分tau}成为

\begin{align}
\int_{0}^{T}\int_{\mathbb{R}^{m}}\left[\partial_{t}\psi\eta(U)+\sum_{\alpha=1}^{m}\partial_{\alpha}\psi Q_{\alpha}(U)\right]\,\mathrm{d}x\mathrm{d}t+\int_{\mathbb{R}^{m}}\psi(x,0)\eta(U_{0}(x))\,\mathrm{d}x\geqslant 0.
\label{align:ch04-熵积分0}
\end{align}

相应地，文献中常作为弱解 $U$ 的熵可容性判据提出的，是\eqref{align:ch04-熵积分0}，而不是稍弱的条件\eqref{align:ch04-熵不等式}。然而应注意，通过\eqref{align:ch04-熵积分0}刻画的可容弱解未必具有所期望的半群性质：即 $U(x,t)$ 可容一般并不蕴含 $U_{\tau}(x,t)=U(x,t+\tau)$ 对所有 $\tau\in[0,T)$ 也可容。因此作者认为，可容性应要么仅通过\eqref{align:ch04-熵不等式}定义，要么通过对所有 $\tau\in[0,T)$ 成立的\eqref{align:ch04-熵积分tau}定义。将来若能证明在某些方程组中 $\mathscr{F}$ 为空，则这一区别便无足轻重。

当熵为凸时，可容弱解所享有增强正则性的第一个迹象，由下述定理给出。

\begin{theorem}\label{theorem:ch04-熵凸正则性}
设 $U(x,t)$ 是\eqref{align:ch04-守恒律组}在 $\mathbb{R}^{m}\times[0,T)$ 上的弱解，它相对于一致凸熵 $\eta$ 满足熵可容性条件\eqref{align:ch04-熵不等式}。则 $t\mapsto U(\cdot,t)$ 在 $[0,T)\setminus\mathscr{F}$ 上于 $L_{\mathrm{loc}}^{p}(\mathbb{R}^{m})$ 中连续，对任一 $p\in[1,\infty)$，其中 $\mathscr{F}$ 至多可数。进而，对某个 $\tau\in[0,T)$，\eqref{align:ch04-熵积分tau}成立当且仅当 $t\mapsto U(\cdot,t)$ 在 $\tau$ 处从右侧于 $L_{\mathrm{loc}}^{p}(\mathbb{R}^{m})$ 中连续，对任一 $p\in[1,\infty)$。
\end{theorem}

\begin{proof}
由于 $t\mapsto U(\cdot,t)$ 与 $t\mapsto\eta(U(\cdot,t))$ 都在 $[0,T)\setminus\mathscr{F}$ 上在 $L^{\infty}(\mathbb{R}^{m})$ weak$^{*}$ 下连续，而 $\eta$ 一致凸，故 $t\mapsto U(\cdot,t)$ 在 $[0,T)\setminus\mathscr{F}$ 上于 $L^{p}(\mathscr{D})$ 中连续，对任一紧子集 $\mathscr{D}\subset\mathbb{R}^{m}$ 与 $p\in[1,\infty)$。

现设对某个 $\tau\in[0,T)$，\eqref{align:ch04-熵积分tau}成立。固定 $\varepsilon>0$，取 $\psi(x,t)=\varphi(x)g(t)$，其中 $\varphi\in C_{0}^{\infty}(\mathbb{R}^{m})$，当 $x\in\mathbb{R}^{m}$ 时 $\varphi(x)\geqslant 0$；而 $g$ 定义为 $g(t)=1-\varepsilon^{-1}(t-\tau)$（当 $0\leqslant t<\tau+\varepsilon$），$g(t)=0$（当 $t+\varepsilon\leqslant t<\infty$）。这给出
\begin{align}
\frac{1}{\varepsilon}\int_{\tau}^{\tau+\varepsilon}\int_{\mathbb{R}^{m}}\varphi(x)[\eta(U(x,\tau))-\eta(U(x,t))]\,\mathrm{d}x\mathrm{d}t\geqslant O(\varepsilon).
\label{align:ch04-熵均值}
\end{align}
由\cref{chapter:ch01-平衡律}，令 $\varepsilon\downarrow 0$，推出 $\operatorname*{ess\,lim}_{t\downarrow\tau}\eta(U(\cdot,t))\leqslant\eta(U(\cdot,\tau))$，其中极限在 $L^{\infty}(\mathscr{D})$ weak$^{*}$ 下取，对任一紧子集 $\mathscr{D}\subset\mathbb{R}^{m}$。回顾 $\operatorname*{lim}_{t\downarrow\tau}U(\cdot,t)=U(\cdot,\tau)$（仍在 $L^{\infty}(\mathscr{D})$ weak$^{*}$ 下），且 $\eta(U)$ 一致凸，故当 $t\downarrow\tau$ 时 $U(\cdot,t)\to U(\cdot,\tau)$，在 $L^{p}(\mathscr{D})$ 中强收敛，对任一 $p\in[1,\infty)$。


反之，设 $t\mapsto U(\cdot,t)$ 在 $\tau$ 处于 $L_{\mathrm{loc}}^{p}(\mathbb{R}^{m})$ 中右连续。固定任一在 $\mathbb{R}^{m}\times[0,T)$ 内紧支集的非负 Lipschitz 函数 $\psi(x,t)$，令 $\phi(x,t)=h(t)\psi(x,t)$，其中 $h(t)=0$（当 $0\leqslant t\leqslant\tau$），$h(t)=\varepsilon^{-1}(t-\tau)$（当 $\tau<t\leqslant\tau+\varepsilon$），$h(t)=1$（当 $\tau+\varepsilon<t<T$）。对检验函数 $\phi$ 应用\eqref{align:ch04-熵不等式}，再令 $\varepsilon\downarrow 0$，即得\eqref{align:ch04-熵积分tau}。证明完成。
\end{proof}

\eqref{align:ch04-熵积分0}或\eqref{align:ch04-熵积分tau}的含义由下述定理进一步阐明。

\begin{theorem}\label{theorem:ch04-总熵衰减}
设 $U(x,t)$ 是\eqref{align:ch04-守恒律组}、\eqref{align:ch04-Cauchy初值}在 $\mathbb{R}^{m}\times[0,T)$ 上的弱解，$U_{0}(\cdot)-\bar{U}\in L^{2}(\mathbb{R}^{m})$，其中 $\bar{U}$ 为某固定状态。设 $U$ 相对于一致凸熵 $\eta$ 满足熵可容性条件\eqref{align:ch04-熵不等式}，并规范化\footnote{如\cref{theorem:ch04-影响区域}证明中所说明。}使 $\eta(\bar{U})=0$，$\mathrm{D}\eta(\bar{U})=0$。则 $U$ 满足\eqref{align:ch04-熵积分0}当且仅当
\begin{align}
\int_{\mathbb{R}^{m}}\eta(U(x,\tau))\,\mathrm{d}x\leqslant\int_{\mathbb{R}^{m}}\eta(U_{0}(x))\,\mathrm{d}x,\qquad 0<\tau<T.
\label{align:ch04-总熵衰减式}
\end{align}
由此可知，\eqref{align:ch04-熵积分tau}对所有 $\tau\in[0,T)$ 成立当且仅当函数 $t\mapsto\int_{\mathbb{R}^{m}}\eta(U(x,t))\,\mathrm{d}x$ 在 $[0,T)$ 上非增。
\end{theorem}

\begin{proof}
设\eqref{align:ch04-熵积分0}成立，固定任一 $\tau\in(0,T)$。对小的 $\varepsilon>0$、$r>0$ 与 $s>0$，令
\begin{align}
h(t)=
\begin{cases}
1&0\leqslant t<\tau\\
\varepsilon^{-1}(\tau-t)+1&\tau\leqslant t<\tau+\varepsilon\\
0&\tau+\varepsilon\leqslant t<T,
\end{cases}
\label{align:ch04-h(t)}
\end{align}
\begin{align}
\phi(x,t)=
\begin{cases}
1&|x|-r-s(\tau-t)<0\\
\varepsilon^{-1}[r+s(\tau-t)-|x|]+1&0\leqslant|x|-r-s(\tau-t)<\varepsilon\\
0&|x|-r-s(\tau-t)\geqslant\varepsilon,
\end{cases}
\label{align:ch04-phi(x,t)}
\end{align}
并对检验函数 $\psi(x,t)=h(t)\phi(x,t)$ 写\eqref{align:ch04-熵积分0}，得
\begin{align}
\frac{1}{\varepsilon}\int_{\tau}^{\tau+\varepsilon}\int_{|x|<r}\eta(U(x,\tau))\,\mathrm{d}x\mathrm{d}t\leqslant\int_{|x|<r+s\tau}\eta(U_{0}(x))\,\mathrm{d}x\nonumber\\
\quad-\frac{1}{\varepsilon}\int_{0}^{\tau}\int_{r+s(\tau-t)<|x|<r+s(\tau-t)+\varepsilon}\left[s\eta(U)+\frac{Q(U)\cdot x}{|x|}\right]\,\mathrm{d}x\mathrm{d}t+O(\varepsilon).
\label{align:ch04-总熵估计}
\end{align}
回顾 $\eta$ 一致凸，$\eta(\bar{U})=0$，$\mathrm{D}\eta(\bar{U})=0$，从而 $\mathrm{D}Q_{\alpha}(\bar{U})=0$，$\alpha=1,\cdots,m$；由 $Q(\bar{U})=0$ 规范化 $Q$，并取 $s$ 充分大，使在解 $U$ 的值域上 $s\eta\geqslant|Q|$ 成立。在\eqref{align:ch04-总熵估计}中令 $\varepsilon\downarrow 0$，用 $\int\eta(U)\,\mathrm{d}x$ 的弱下半连续性，再令 $r\uparrow\infty$，即得\eqref{align:ch04-总熵衰减式}。

反之，设\eqref{align:ch04-总熵衰减式}成立。由于当 $\tau\downarrow 0$ 时 $U(\cdot,\tau)\to U_{0}(\cdot)$ 在 $L^{\infty}(\mathbb{R}^{m})$ weak$^{*}$ 下，\eqref{align:ch04-总熵衰减式}与 $\int\eta(V)\,\mathrm{d}x$ 的弱下半连续性蕴含 $U(\cdot,\tau)\to U_{0}(\cdot)$ 在 $L_{\mathrm{loc}}^{p}(\mathbb{R}^{m})$ 中。再由\cref{theorem:ch04-熵凸正则性}，这给出\eqref{align:ch04-熵积分0}。证明完成。
\end{proof}

\eqref{align:ch04-总熵衰减式}所体现的总熵单调衰减在物理上是吸引人的。它基于可容性条件\eqref{align:ch04-熵积分tau}，并配合映射 $V\mapsto\int\eta(V)\,\mathrm{d}x$ 的弱下半连续性；只要 $V$ 不受限制，这一半连续性对 $\eta(V)$ 的凸性既是必要的又是充分的。然而，在第 III 章中我们遇到力学与电动力学中出现的重要方程组，其中熵函数不允许为凸。救场的是：状态向量须满足附加条件，即使 $\eta(V)$ 不凸，也可能使 $\int\eta(V)\,\mathrm{d}x$ 弱下半连续。

作为说明，考虑支配热弹性介质等熵运动的方程组\cref{chapter:ch03-双曲平衡律方程组}。回顾熵函数 $\eta=\varepsilon(F)+\tfrac{1}{2}|v|^{2}$ 之所以不能为凸，是因为内能函数 $\varepsilon(F)$ 不允许为凸。尽管如此，已知（文献见第 4.9 小节），由于 $F$ 被约束为梯度，映射 $F\mapsto\int\varepsilon(F)\,\mathrm{d}x$ 在 $L^{\infty}$ weak$^{*}$ 下下半连续当且仅当 $\varepsilon(F)$ 在 Morrey 意义下是\emph{拟凸}的，即记 $\mathscr{K}$ 为 $\mathbb{R}^{3}$ 中的标准超立方体，对每一常矩阵 $\hat{F}$ 与任一在 $\mathscr{K}$ 内紧支集的 Lipschitz 向量场 $\chi$，

\begin{align}
\varepsilon(\hat{F})\leqslant\int_{\mathscr{K}}\varepsilon(\hat{F}+\nabla\chi)\,\mathrm{d}x
\label{align:ch04-拟凸}
\end{align}
成立。用物理语言，\eqref{align:ch04-拟凸}规定：当 $\mathscr{K}$ 的任一均匀形变在所有具有相同边界值的 $\mathscr{K}$ 放置中使贮于 $\mathscr{K}$ 内的内能最小时，$\varepsilon$ 是拟凸的。

$\varepsilon(F)$ 的凸性是拟凸的充分条件，但非必要条件。另一方面，$s$ 为常数的秩一凸性\cref{chapter:ch03-双曲平衡律方程组}是 $\varepsilon(F)$ 拟凸的必要条件，但非充分条件。事实上，由于\eqref{align:ch04-拟凸}是非局部的，要检验任一特定的秩一凸但非凸函数是否拟凸并不容易。


构造物理可容、非凸但拟凸的内能的一种方法基于：当 $\varepsilon(F)=\varphi(F)$ 时\eqref{align:ch04-拟凸}取等式，其中 $\varphi(F)$ 是形如\cref{chapter:ch02-连续介质物理导引}的任一零 Lagrangian。于是零 Lagrangian 在 $L^{\infty}$ weak$^{*}$ 下连续。由此，本构方程形如

\begin{align}
\varepsilon(F)=\theta(F,F^{*},\det F)
\label{align:ch04-多凸}
\end{align}
的内能（其中 $\theta(F,H,\delta)$ 在 $\mathbb{R}^{19}$ 中为凸）在 $L^{\infty}$ weak$^{*}$ 下下半连续，从而拟凸。这种类型的本构方程称为\emph{多凸}的。它们为实际弹性材料提供了现实模型，并在弹性静力学中起重要作用。它们在弹性动力学中的作用将在\cref{chapter:ch05-熵与经典解稳定性}中阐明。就当前目的而言，当内能函数为多凸时，\cref{theorem:ch04-总熵衰减}的断言对等温热弹性方程组\cref{chapter:ch03-双曲平衡律方程组}成立。


正如我们稍后将看到，熵可容性条件排除了双曲守恒律组不想要的、赝的弱解中的一部分，但未必是全部。一个可能的补救是要求\eqref{align:ch04-熵积分0}对系统的每一凸熵同时成立。然而，这似乎仅对配有丰富熵族的极特殊方程组类有前途，特别是正如我们将在第 VI 章中看到的对标量守恒律。在这一联系中，应注意以下注记。

\begin{remark}\label{remark:ch04-多熵可容性}
设\eqref{align:ch04-守恒律组}配有一个熵--熵流对 $(\bar{\eta},\bar{Q})$，且 $\mathrm{D}^{2}\bar{\eta}$ 在 $\mathcal{O}$ 上正定；并对\eqref{align:ch04-守恒律组}的解 $U$ 考虑更强的可容性条件，即\eqref{align:ch04-熵积分0}须对每一凸熵的熵--熵流对 $(\eta,Q)$ 成立。此时，对任一熵--熵流对 $(\eta,Q)$，$\partial_{t}\eta(U)+\operatorname{div}Q(U)$ 将是一个测度。这由下述观察推出：任一 $C^{2}$ 熵--熵流对 $(\eta,Q)$ 可写成熵--熵流对 $(k\bar{\eta}+\eta,k\bar{Q}+Q)$ 与 $(k\bar{\eta},k\bar{Q})$ 之差；当 $k$ 充分大时，$k\bar{\eta}+\eta$ 与 $k\bar{\eta}$ 的 Hessian 在 $\mathcal{O}$ 的任一紧子集上都正定。
\end{remark}

在缺乏丰富熵族、且对单一熵的\eqref{align:ch04-熵不等式}不足以排除所有赝解时，可尝试采取一种更具选择性的可容性条件，以挑出物理可容解，这始终符合热力学第二定律的精神。由于\eqref{align:ch04-总熵衰减式}表明可容解必为耗散的，自然要问：以最高速率耗散的解是否具有特殊地位。人们可尝试各种最大耗散率的刻画。于是，在\cref{theorem:ch04-总熵衰减}的设定下，解 $U$ 可称为\emph{极大耗散}的，如果对每一 $t\in[0,T)$ 以及在 $\mathbb{R}^{m}\times[0,t]$ 上与 $U$ 重合的任一其他解 $\hat{U}$，下列二者之一成立：
\begin{align}
D^{+}\int_{\mathbb{R}^{m}}\eta(U(x,t))\,\mathrm{d}x\leqslant D_{+}\int_{\mathbb{R}^{m}}\eta(\hat{U}(x,t))\,\mathrm{d}x,
\label{align:ch04-极大耗散}
\end{align}
其中 $D^{+}$ 与 $D_{+}$ 分别表示上、下右 Dini 导数；或者，等价地，存在递减序列 $\{t_{k}\}$，$t_{k}\to t$（当 $k\to\infty$），使得
\begin{align}
\int_{\mathbb{R}^{m}}\eta(U(x,t_{k}))\,\mathrm{d}x\leqslant\int_{\mathbb{R}^{m}}\eta(\hat{U}(x,t_{k}))\,\mathrm{d}x,\qquad k=1,2,\cdots.
\label{align:ch04-极大耗散序列}
\end{align}
注意\eqref{align:ch04-极大耗散}与\eqref{align:ch04-极大耗散序列}精神相同，但互不蕴含。我们将在第 4.8 小节中简略、并在第 9.7 小节中更详尽地回到这些考虑。

当可容解 $U$ 属于 $BV_{\mathrm{loc}}$ 类时，\cref{chapter:ch01-平衡律}蕴含熵产生测度集中于 $U$ 的近似跳跃间断点集上，即激波面上。此时，\eqref{align:ch04-熵不等式}化为局部条件\cref{chapter:ch01-平衡律}，在当前记号下取形式

\begin{align}
-s[\eta(U_{+})-\eta(U_{-})]+[Q(U_{+})-Q(U_{-})]\nu\leqslant 0.
\label{align:ch04-激波熵条件}
\end{align}
要判定 $U$ 相对于熵 $\eta$ 的可容性，须在沿方向 $\nu\in\mathbb{S}^{m-1}$ 以速度 $s$ 传播的激波的任一点处检验\eqref{align:ch04-激波熵条件}。

作为应用，让我们相对于熵--熵流对 $(\tfrac{1}{2}u^{2},\tfrac{1}{3}u^{3})$ 检验由\eqref{align:ch04-非唯一解族}定义的解族 $u_{\alpha}(x,t)$ 对问题\eqref{align:ch04-Burgers}、\eqref{align:ch04-Riemann初值}的可容性。显然，对任一 $\alpha\in(0,1]$，静止激波 $x=0$ 违反\eqref{align:ch04-激波熵条件}。因而该族中唯一可容的解是 $u_{0}(x,t)$，它在远离原点处 Lipschitz 连续。

\end{document}
